\documentclass[8pt,aspectratio=169]{beamer}
\usetheme[progressbar=frametitle, block=fill]{metropolis}

\usepackage{fontspec}
\setsansfont{FiraSans}[
  Path=/usr/local/texlive/2026/texmf-dist/fonts/opentype/public/fira/,
  Extension=.otf,
  UprightFont=*-Light,
  BoldFont=*-Regular,
  ItalicFont=*-LightItalic,
  BoldItalicFont=*-Italic
]
\setmainfont{FiraSans}[
  Path=/usr/local/texlive/2026/texmf-dist/fonts/opentype/public/fira/,
  Extension=.otf,
  UprightFont=*-Light,
  BoldFont=*-Regular,
  ItalicFont=*-LightItalic,
  BoldItalicFont=*-Italic
]
\setmonofont{FiraMono}[
  Path=/usr/local/texlive/2026/texmf-dist/fonts/opentype/public/fira/,
  Extension=.otf,
  UprightFont=*-Regular,
  BoldFont=*-Medium,
  ItalicFont=*-Oblique,
  BoldItalicFont=*-MediumOblique
]
\usepackage[vietnamese]{babel}
\usepackage{amsmath,amssymb,amsfonts,mathtools,bm}
\usepackage{graphicx,booktabs,array,multirow,colortbl}
\usepackage{tikz,pgfplots}
\pgfplotsset{compat=1.18}
\usetikzlibrary{shapes.geometric,arrows.meta,positioning,fit,backgrounds,calc,decorations.pathreplacing}
\usepackage{xcolor,hyperref,appendixnumberbeamer,microtype}

\graphicspath{{figures/}{../../results/runs_all/figures/two_moons/}{../../results/runs_all/figures/two_moons/deep/}}

\definecolor{deepblue}{HTML}{1B3A6B}
\definecolor{skyblue}{HTML}{3A7CA5}
\definecolor{accent}{HTML}{F4A261}
\definecolor{teal}{HTML}{2A9D8F}
\definecolor{violet}{HTML}{5E548E}
\definecolor{lightgray}{HTML}{F0F0F0}
\definecolor{proposegreen}{HTML}{1B6B3A}
\definecolor{highlight}{HTML}{C0392B}

\setbeamercolor{frametitle}{bg=deepblue}
\setbeamercolor{progress bar}{fg=accent,bg=skyblue!30}
\setbeamercolor{alerted text}{fg=accent}
\setbeamercolor{block title}{bg=skyblue!20,fg=deepblue}
\setbeamercolor{block body}{bg=lightgray}
\setbeamercolor{block title example}{bg=teal!20,fg=teal!80!black}
\setbeamercolor{block body example}{bg=teal!5}
\setbeamercolor{block title alerted}{bg=accent!25,fg=deepblue}
\setbeamercolor{block body alerted}{bg=accent!8}
\setbeamercolor{palette primary}{bg=deepblue}
\setbeamerfont{title}{size=\Large,series=\bfseries}
\setbeamerfont{frametitle}{size=\normalsize,series=\bfseries}
\metroset{sectionpage=progressbar}

\newcommand{\hl}[1]{\textcolor{accent}{\textbf{#1}}}
\newcommand{\tc}[1]{\textcolor{teal}{#1}}
\newcommand{\bad}[1]{\textcolor{highlight}{\textbf{#1}}}
\newcommand{\good}[1]{\textcolor{proposegreen}{\textbf{#1}}}
\newcommand{\opt}[1]{\texttt{#1}}

% --- F9 optimizer colours (khớp OPT_COLORS trong verification_common.py) ---
\definecolor{optsgd}{HTML}{1f77b4}
\definecolor{optsgdm}{HTML}{17becf}
\definecolor{optadam}{HTML}{9467bd}
\definecolor{optadah}{HTML}{d62728}
\definecolor{optsassha}{HTML}{ff7f0e}
\definecolor{optsophia}{HTML}{8c564b}
\definecolor{optkfac}{HTML}{2ca02c}
% một mục legend: màu, kiểu nét (dashed/solid), nhãn
\newcommand{\optline}[3]{\tikz[baseline=-0.5ex]{\draw[#1,line width=1.4pt,#2](0,0)--(0.45,0);}~#3}
% dải legend đầy đủ cho hình F9 (nét đứt = 1st-order, nét liền = 2nd-order)
\newcommand{\fninelegend}{%
  \footnotesize
  \optline{optsgd}{dashed}{SGD}\enspace
  \optline{optsgdm}{dashed}{SGD-mom}\enspace
  \optline{optadam}{dashed}{Adam}\enspace
  \optline{optadah}{solid}{AdaHessian}\enspace
  \optline{optsassha}{solid}{SASSHA}\enspace
  \optline{optsophia}{solid}{SophiaH}\enspace
  \optline{optkfac}{solid}{KFAC}%
  \\[1pt]{\scriptsize nét đứt = 1st-order \quad nét liền = 2nd-order}%
}

\title{Optimizer có giúp làm thuyên giảm\\Loss of Plasticity?}
\subtitle{Mức độ \& cơ chế --- soi trên two-moons\\Báo cáo tiến độ nghiên cứu}
\author{Nguyễn Lâm Phú Quý, Bàng Mỹ Linh}
\date{\today}

\begin{document}

\begin{frame}[plain]
  \titlepage
\end{frame}

\begin{frame}{Nội dung trình bày}
  \setbeamertemplate{section in toc}[sections numbered]
  \tableofcontents[hideallsubsections]
\end{frame}

\section{Bối cảnh \& câu hỏi gốc}

\begin{frame}{Loss of Plasticity là gì?}
\vspace{1.0em}
\begin{block}{Định nghĩa}
\large
Trong học liên tục, mạng có thể \textbf{không chỉ quên task cũ}, mà còn
\hl{mất dần khả năng học task mới}. Hiện tượng này gọi là
\textit{Loss of Plasticity} (LoP).
\end{block}

\vspace{1.0em}
\begin{columns}[T]
\column{0.48\textwidth}
\begin{alertblock}{Không phải câu hỏi forgetting}
\small
Không hỏi ``còn nhớ task cũ không?''; tín hiệu chính là:
\[
\text{task mới tới} \quad \Rightarrow \quad
\text{còn fit được task hiện tại không?}
\]
\end{alertblock}

\column{0.48\textwidth}
\begin{exampleblock}{Góc nhìn cơ chế}
\small
Gradient descent có thể rơi vào các \textbf{invariant submanifold}: một khi vào đó, gradient gần như chỉ chạy bên trong bẫy, nên capacity hữu ích cho task mới bị khóa lại.
\end{exampleblock}
\end{columns}
\end{frame}

\begin{frame}{Các dấu hiệu cơ chế cần đo}
\small
\begin{columns}[T]
\column{0.48\textwidth}
\begin{alertblock}{Dormant / frozen units}
Unit gần như không kích hoạt trên dữ liệu hiện tại; gradient về trọng số vào biến mất $\Rightarrow$ capacity bị đóng băng.
\end{alertblock}

\begin{alertblock}{Cloned / duplicate units}
Các unit dư thừa có activation giống nhau. Theo ajoudaki-lop, nếu chúng nhận forward/backward giống nhau thì gradient cũng giống nhau, nên training có thể giữ nguyên cấu trúc clone.
\end{alertblock}

\begin{alertblock}{Rank collapse}
Effective rank / stable rank của activation giảm $\Rightarrow$ biểu diễn mất đa dạng, nhiều unit nhưng ít chiều hữu ích.
\end{alertblock}

\column{0.48\textwidth}
\begin{alertblock}{Weight explosion}
$\|W\|$ tăng mạnh $\Rightarrow$ activation dễ trôi vào vùng chết / bão hòa.
\end{alertblock}

\begin{alertblock}{Hessian spectral collapse}
Phổ curvature co hẹp hoặc suy biến $\Rightarrow$ mất các hướng cong hữu ích để học task mới.
\end{alertblock}

\begin{exampleblock}{Điểm chung}
Các metric này không phải kết luận riêng lẻ; chúng là các cửa sổ để xem optimizer đang đẩy mạng vào bẫy nào.
\end{exampleblock}
\end{columns}
\end{frame}

\begin{frame}{Optimizer bậc 2 đã chạy rộng trên nhiều benchmark}
\small
\begin{block}{Bối cảnh}
Hướng \textbf{optimizer bậc 2} không chỉ có một curve đẹp: bộ thí nghiệm trải rộng trên \textbf{MLP, ConvNet, ResNet, PPO}, với \textbf{CBP} làm baseline giữ plasticity và nhiều optimizer curvature-aware kết hợp SDP.
\end{block}

\vspace{0.35em}
\centering
\renewcommand{\arraystretch}{1.14}
\resizebox{0.96\linewidth}{!}{%
\begin{tabular}{p{0.22\linewidth}p{0.30\linewidth}p{0.27\linewidth}p{0.22\linewidth}}
\toprule
Benchmark & Quy mô đã chạy & Các phương pháp so sánh & Tín hiệu mạnh nhất \\
\midrule
Permuted-MNIST & 800 tasks, MLP 2000$\times$2, no BN, 5 seeds
& CBP; K-FAC (NGD); AdaHessian/SophiaH/SASSHA $\pm$ SDP
& K-FAC (no SDP): 0.00 dormant, rank 1733.6; AdaHessian+SDP: 97.04 acc \\
ImageNet32 binary & 2000 binary tasks, ConvNet, no BN, 5 seeds
& CBP; AdaHessian/SophiaH/SASSHA $\pm$ SDP
& SASSHA+SDP: 90.09 acc, 0.00 dormant \\
CIFAR-100 incremental & 20 tasks $\times$ 5 classes, ResNet-18+BN, 5 seeds
& SGD/Adam, CBP, K-FAC (NGD), Hessian methods, SDP
& SASSHA+SDP: 71.35 acc, 0.20 dormant \\
Ant-v4 PPO & 50M env steps, policy MLP, 5 seeds
& AdaHessian+SDP, SASSHA+SDP, ACKTR/OPC ref.
& AdaHessian+SDP: return 3361.9; SASSHA+SDP: 1.12 dormant \\
\bottomrule
\end{tabular}}

\begin{exampleblock}{Ý chính}
Optimizer bậc 2 đã được kiểm chứng trên nhiều regime LoP thực; phần two-moons tiếp theo là bước \textbf{soi cơ chế} cho riêng câu hỏi optimizer-only.
\end{exampleblock}
\end{frame}

\begin{frame}{So với CBP: curvature-aware + SDP có thể vượt baseline mạnh}
\small
\centering
\renewcommand{\arraystretch}{1.14}
\resizebox{0.92\linewidth}{!}{%
\begin{tabular}{llccc}
\toprule
Benchmark & Method & Final Acc.~$\uparrow$ & Dormant~$\downarrow$ & Stable rank~$\uparrow$ \\
\midrule
\multirow{3}{*}{Permuted-MNIST}
& CBP & $95.77_{\pm0.12}$ & $0.04_{\pm0.01}$ & $1473_{\pm42}$ \\
& \good{K-FAC (NGD), no SDP} & $95.70$ & $\mathbf{0.00}$ & $\mathbf{1733.6}$ \\
& \good{AdaHessian+SDP} & $\mathbf{97.04}_{\pm0.14}$ & $1.41_{\pm0.18}$ & $1674_{\pm39}$ \\
\midrule
\multirow{2}{*}{ImageNet32 binary}
& CBP & $89.62_{\pm0.31}$ & $0.02_{\pm0.01}$ & $\mathbf{118.81}_{\pm3.2}$ \\
& \good{SASSHA+SDP} & $\mathbf{90.09}_{\pm0.28}$ & $\mathbf{0.00}_{\pm0.00}$ & $112.84_{\pm2.6}$ \\
\bottomrule
\end{tabular}}

\vspace{0.2em}
{\scriptsize $\uparrow$ cao hơn = tốt hơn;\quad $\downarrow$ thấp hơn = tốt hơn.}

\vspace{0.45em}
\begin{columns}[T]
\column{0.48\textwidth}
\begin{alertblock}{Vượt CBP ở accuracy}
\small
MNIST: +1.27 điểm acc và rank cao hơn.\\
ImageNet32: +0.47 điểm acc, dormant về 0.00.
\end{alertblock}
\column{0.48\textwidth}
\begin{block}{K-FAC: tín hiệu optimizer-only}
\small
Đáng chú ý: \good{K-FAC không cần SDP} đã đạt dormant \textbf{0.00} và rank cao nhất (1733.6), ngang CBP — gợi ý chính \textbf{preconditioner} có thể tự giữ plasticity, đúng câu hỏi ta đang theo.
\end{block}
\end{columns}
\end{frame}

\begin{frame}{Kết quả bổ sung: CIFAR-100 và RL cũng có tín hiệu mạnh}
\small
\centering
\renewcommand{\arraystretch}{1.13}
\resizebox{0.94\linewidth}{!}{%
\begin{tabular}{llccc}
\toprule
Setting & Method & Performance~$\uparrow$ & Dormant~$\downarrow$ & Ghi chú \\
\midrule
\multirow{3}{*}{CIFAR-100, ResNet-18+BN}
& AdaHessian & $62.4_{\pm1.1}$ acc & $4.1_{\pm0.4}$ & baseline Hessian \\
& \good{AdaHessian+SDP} & $71.2_{\pm0.8}$ acc & $1.8_{\pm0.3}$ & +8.8 điểm acc \\
& \good{SASSHA+SDP} & $\mathbf{71.35}_{\pm0.7}$ acc & $\mathbf{0.20}$ & dormant gần 0 \\
\midrule
\multirow{3}{*}{Ant-v4 PPO, 50M steps}
& \good{AdaHessian+SDP} & $\mathbf{3361.9}_{\pm184}$ return & $3.50_{\pm0.42}$ & rank $237.1_{\pm8.2}$ \\
& \good{SASSHA+SDP} & $2939.9_{\pm212}$ return & $\mathbf{1.12}_{\pm0.18}$ & rank $235.1_{\pm9.1}$ \\
& ACKTR + OPC & $4135.5$ return & $\mathbf{0.0}$ & tham chiếu (K-FAC RL) \\
\bottomrule
\end{tabular}}

\vspace{0.2em}
{\scriptsize $\uparrow$ cao hơn = tốt hơn;\quad $\downarrow$ thấp hơn = tốt hơn.}

\vspace{0.35em}
\begin{alertblock}{Nhưng đây vẫn là kết quả có can thiệp}
\small
SDP tái định hình phổ trọng số ở task boundary; CBP thì thay neuron low-utility. Vì vậy câu hỏi còn mở là: \textbf{nếu bỏ mọi can thiệp, bản thân optimizer/preconditioner làm gì với cơ chế LoP?}
\end{alertblock}
\end{frame}

\begin{frame}{Từ benchmark sang phân tích cơ chế}
\vspace{0.2em}
\begin{alertblock}{Khả quan, nhưng bức tranh quá phân tán}
\small
Tín hiệu tốt nhưng trải mỏng trên \hl{nhiều benchmark $\times$ nhiều optimizer $\times$ nhiều can thiệp} --- mỗi ô là một con số riêng, khó gộp thành một câu chuyện thống nhất để \textbf{nghiên cứu và viết bài}.
\end{alertblock}

\vspace{0.6em}
\centering
\begin{tikzpicture}[
  font=\footnotesize,
  scatter/.style={rectangle, rounded corners=3pt, draw=highlight, fill=highlight!8,
    text width=4.1cm, align=center, inner sep=6pt},
  unified/.style={rectangle, rounded corners=3pt, draw=proposegreen, fill=proposegreen!10,
    text width=4.1cm, align=center, inner sep=6pt},
  flow/.style={-{Stealth[length=3mm]}, line width=1pt, accent}
]
\node[scatter] (a) {\bad{4 benchmark $\times$ 7 optimizer}\\$\times$ \{SDP, OPC, CBP, $\gamma$\}\\[2pt] \textit{bảng số rời rạc}};
\node[unified, right=2.4cm of a] (b) {\good{Một khung cơ chế chung}\\dormant · rank · sharpness · route\\[2pt] \textit{giải thích thống nhất}};
\draw[flow] (a) -- node[above=1pt]{\scriptsize chuyển góc nhìn} (b);
\end{tikzpicture}

\vspace{0.6em}
\begin{exampleblock}{Vì sao chọn mechanism analysis}
\small
\begin{itemize}\setlength\itemsep{1pt}
\item \tc{Tự nhiên hơn}: trả lời \emph{tại sao} LoP xảy ra, thay vì chỉ xếp hạng con số.
\item \tc{Tổng quát hơn}: một lời giải áp cho mọi optimizer, ít lệ thuộc benchmark/can thiệp cụ thể.
\item \tc{Đo được mức độ \& cơ chế}: đúng câu hỏi ``optimizer giúp ở mức nào, nhờ cơ chế nào''.
\end{itemize}
\end{exampleblock}
\end{frame}

\begin{frame}{Câu hỏi nghiên cứu}
\vspace{0.4em}
\begin{center}
\fcolorbox{accent}{accent!10}{%
\begin{minipage}{0.84\textwidth}
\centering
\vspace{0.5em}
{\Large\textbf{Optimizer --- chỉ qua preconditioner}}\\[0.25em]
{\Large$\Delta\theta=-\eta P^{-1}\nabla_\theta\mathcal L$}\\[0.35em]
{\large có giúp \hl{làm thuyên giảm} LoP không?}\\[0.5em]
\textbf{(1) Ở mức độ nào? \qquad (2) Nhờ cơ chế nào?}
\vspace{0.5em}
\end{minipage}}
\end{center}

\vspace{1em}
\begin{columns}[c]
\column{0.43\textwidth}
\centering
\bad{1st-order}\\[0.3em]
$\Delta\theta=-\eta\nabla\mathcal L$\\
\scriptsize Gradient-only, không dùng curvature trực tiếp.
\column{0.10\textwidth}
\centering
{\Large vs.}
\column{0.43\textwidth}
\centering
\good{Curvature-aware}\\[0.3em]
$\Delta\theta=-\eta P^{-1}\nabla\mathcal L$\\
\scriptsize Precondition bằng Hessian / Fisher / sharpness.
\end{columns}

\vspace{0.8em}
\begin{exampleblock}{Giả thuyết ngây thơ H}
\textit{1st-order kẹt trong sub-manifold cũ; curvature-aware optimizer thoát sang basin mới.}
\end{exampleblock}
\end{frame}

\section{Testbed two-moons}

\begin{frame}{Vì sao two-moons \& thiết kế thí nghiệm}
\begin{columns}[T]
\column{0.50\textwidth}
\begin{block}{Mục tiêu}
Tách bạch cơ chế optimizer khỏi can thiệp: \textbf{không SDP, không EMA, không CBP, không shrink-and-perturb}.
\end{block}

\begin{itemize}\small
\item Stream: \hl{200 task}, xoay mượt \hl{1.8$^\circ$/task}, không revisit.
\item Mạng: MLP ReLU \hl{512$\times$4}, 3 epoch/task.
\item Factorial: 7 optimizer $\times$ 2 BatchNorm $\times$ \hl{seed 42}.
\item Optimizer: SGD, SGD-mom, Adam, AdaHessian, SophiaH, SASSHA, KFAC.
\end{itemize}

\begin{block}{Quy ước ký hiệu}
\small \textbf{BN=0} = \textbf{không} BatchNorm; \quad \textbf{BN=1} = \textbf{có} BatchNorm.
\end{block}

\column{0.47\textwidth}
\begin{exampleblock}{Điểm mấu chốt}
Log đủ \texttt{theta\_points}: $201\times790{,}530$ tham số. Từ đó tái dựng bằng numpy từng unit trong 2048 hidden units; sai số validate dormant: \hl{0.0000}.
\end{exampleblock}

\begin{block}{6 cơ chế đo}
\small dormant, effective rank, NTK condition number, trajectory curvature, avg $\|W\|$, participation ratio.
\end{block}
\end{columns}
\end{frame}

\section{Phán quyết chính}

\begin{frame}{F9: zoom theo từng nhóm cơ chế}
\begin{columns}[T]
\column{0.49\textwidth}
\begin{block}{Convention}
\small
\textbf{Nét đứt} = 1st-order; \textbf{nét liền} = curvature-aware.\\[0.35em]
Mỗi nhóm sau phóng to 2 hàng cơ chế, giữ hai cột BN=0 / BN=1 (\textbf{BN=0}=không BatchNorm, \textbf{BN=1}=có BatchNorm).
\end{block}

\column{0.49\textwidth}
\begin{alertblock}{Takeaway trước khi zoom}
\small
Nếu H đúng theo trục 1st/2nd, nét đứt và nét liền phải tách rõ. Nhưng tín hiệu nổi bật lại là \bad{Adam} + \bad{SophiaH}: một 1st-order, một 2nd-order clamped.
\end{alertblock}
\end{columns}

\vspace{0.6em}
\begin{columns}[T]
\column{0.52\textwidth}
\centering
\small
\renewcommand{\arraystretch}{1.25}
\begin{tabular}{lll}
\toprule
Frame & Nhóm cơ chế & Ý nghĩa khi zoom \\
\midrule
F9-A & rank + frozen units & Capacity có sụp? \\
F9-B & NTK + trajectory & Hình học tối ưu suy biến? \\
F9-C & weight + participation & Norm/subspace khóa lại? \\
\bottomrule
\end{tabular}

\column{0.46\textwidth}
\begin{block}{Tinh chỉnh câu hỏi nghiên cứu}
\small
$\Rightarrow$ Câu hỏi nên đặt \textbf{chi tiết hơn}: khác biệt về LoP có thể đến từ \textbf{dạng preconditioner} của optimizer, chứ không phải ranh giới 1st/2nd-order.
\end{block}
\end{columns}
\end{frame}

\begin{frame}{F9-A: Capacity collapse --- rank + frozen units}
\centering
\includegraphics[width=\linewidth,trim=0 733bp 0 45bp,clip]{F9_lop_mechanism.png}

\vspace{0.15em}{\centering\fninelegend\par}\vspace{0.25em}
\begin{alertblock}{Tín hiệu capacity}
\small
Ở BN=0, \bad{SophiaH} tăng frozen rất sớm và effective-rank rơi xuống đáy; \bad{Adam} cũng là death/rank-collapse route. Sang BN=1, đa số optimizer được cứu khỏi frozen, nhưng Adam vẫn tụt rank và dormant leo cao. Capacity collapse vì vậy không đi theo nét đứt/liền.
\end{alertblock}
\end{frame}

\begin{frame}{F9-B: Degeneracy signals --- NTK + trajectory}
\centering
\includegraphics[width=\linewidth,trim=0 401bp 0 395bp,clip]{F9_lop_mechanism.png}

\vspace{0.15em}{\centering\fninelegend\par}\vspace{0.25em}
\begin{alertblock}{Tín hiệu hình học tối ưu}
\small
NTK condition đẩy \bad{Adam}/\bad{SophiaH} lên nhiều bậc độ lớn, cùng hướng với frozen/rank collapse. Trajectory curvature chỉ là tín hiệu phụ: stream two-moons tự xoay nên đường đi cong là tự nhiên, không đủ để gán nguyên nhân LoP.
\end{alertblock}
\end{frame}

\begin{frame}{F9-C: Norm/subspace --- weight + participation}
\centering
\includegraphics[width=\linewidth,trim=0 69bp 0 695bp,clip]{F9_lop_mechanism.png}

\vspace{0.15em}{\centering\fninelegend\par}\vspace{0.25em}
\begin{alertblock}{Tín hiệu norm/subspace}
\small
Norm cho thấy \bad{SophiaH} là ca blow-up rõ nhất, không phải toàn bộ nhóm 1st-order. Participation ratio lại đặt SGD ở ít chiều nhất, hơi hợp giả thuyết subspace-trap ban đầu, nhưng lành tính trong regime này vì SGD vẫn giữ accuracy.
\end{alertblock}
\end{frame}

\begin{frame}{Từ F9: vì sao trục 1st/2nd chưa đủ}
\small
\begin{columns}[T]
\column{0.56\textwidth}
\renewcommand{\arraystretch}{1.15}
\resizebox{\linewidth}{!}{%
\begin{tabular}{p{0.26\linewidth}p{0.66\linewidth}}
\toprule
Cơ chế & Quan sát sơ bộ trên two-moons \\
\midrule
Dormant & BN=0: SophiaH, Adam cao vọt; optimizer còn lại thấp. \\
Eff. rank & SophiaH / Adam sụp mạnh; SGD-mom sụp qua route khác. \\
NTK cond. & Adam + SophiaH suy biến rõ nhất. \\
$\|W\|$ & Blow-up nằm ở SophiaH, không phải toàn bộ 1st-order. \\
Participation ratio & SGD ít chiều nhất, nhưng không mất accuracy. \\
\bottomrule
\end{tabular}}

\column{0.41\textwidth}
\begin{alertblock}{Tín hiệu thay thế, chưa phải kết luận}
Một trục đáng nghi:
\[
\bad{\text{diagonal/clamped adaptive}}
\quad \text{vs.} \quad
\good{\text{full curvature-aware}}.
\]
Adam và SophiaH cùng rơi vào nhóm đầu, dù khác ``bậc'' optimizer.
\end{alertblock}

\begin{exampleblock}{Cần kiểm chứng}
\centering
\textbf{$\geq$5 seed + regime đột ngột} trước khi phát biểu lại thành claim chính.
\end{exampleblock}
\end{columns}
\end{frame}

\begin{frame}{Dormancy cao không chỉ có một nghĩa: churn vs frozen-trap}
\centering
\small
\renewcommand{\arraystretch}{1.15}
\resizebox{0.88\linewidth}{!}{%
\begin{tabular}{lcccc}
\toprule
Optimizer (BN=0) & \%dead TB & $P(\text{hồi sinh}|\text{chết})$ & $P(\text{chết}|\text{sống})$ & \%chết vĩnh viễn \\
\midrule
sgd & 9.3 & 8.5 & 0.9 & 2.7 \\
sgd\_momentum & 9.2 & 6.4 & 0.7 & 2.0 \\
\bad{adam} & \bad{56.3} & \bad{13.6} & \bad{18.0} & \textbf{0.0} \\
adahessian & 10.3 & 8.6 & 1.0 & 2.4 \\
\bad{sophiah} & \bad{81.5} & \bad{2.1} & 9.9 & \bad{15.6} \\
\good{sassha} & 7.4 & \good{12.9} & 1.0 & \good{1.2} \\
kfac & 9.4 & 8.7 & 0.9 & 4.0 \\
\bottomrule
\end{tabular}}

\vspace{0.4em}
\begin{alertblock}{Hai kiểu dormancy khác nhau}
\bad{Adam} = churn hỗn loạn: chết nhiều nhưng hồi sinh/đổi trạng thái mạnh. \bad{SophiaH} = frozen-trap thật: hồi sinh yếu và có 15.6\% unit chết xuyên suốt.
\end{alertblock}
\end{frame}

\section{Đào sâu cơ chế}

\begin{frame}{Ba route chuyển trạng thái: death, cloning, diversifying}
\begin{columns}[T]
\column{0.55\textwidth}
\centering
\includegraphics[width=\linewidth,height=0.68\textheight,keepaspectratio]{F16_state_routes.png}
\column{0.42\textwidth}
\small
\renewcommand{\arraystretch}{1.1}
\resizebox{\linewidth}{!}{%
\begin{tabular}{lrrrl}
\toprule
Opt & $\Delta$dead & $\Delta$clone & $\Delta$unique & Route \\
\midrule
sgd & -0.3 & -4.9 & +5.1 & \good{diversify} \\
adahessian & -0.7 & -2.2 & +2.9 & \good{diversify} \\
sassha & -2.3 & -0.3 & +2.6 & \good{diversify} \\
kfac & -1.4 & -3.7 & +5.0 & \good{diversify} \\
sgd-mom & +1.4 & +5.8 & -7.2 & cloning \\
\bad{adam} & +48.1 & -43.7 & -4.5 & \bad{death} \\
\bad{sophiah} & +39.4 & -36.2 & -3.3 & \bad{death} \\
\bottomrule
\end{tabular}}

\begin{alertblock}{Hướng chuyển trạng thái mới là tín hiệu}
Cloning là trạng thái khởi đầu chung do dữ liệu 2-D; cái phân biệt là \textbf{hướng chuyển trạng thái}. Nhóm khỏe đi từ cloned sang unique, SGD-mom tăng clone, còn Adam/SophiaH chuyển capacity sang dead.
\end{alertblock}
\end{columns}
\end{frame}

\begin{frame}{Structural vs flicker death: lõi chết có thật}
\begin{columns}[T]
\column{0.55\textwidth}
\centering
\includegraphics[width=\linewidth,height=0.68\textheight,keepaspectratio]{F19_structural_vs_flicker.png}
\column{0.42\textwidth}
\small
\begin{block}{Cách đo}
Cố định $\theta_{\mathrm{final}}$, quét dead trên đủ 200 hướng dữ liệu. \textbf{Structural dead} = chết trên mọi hướng; flicker = chết ở hướng cuối nhưng sống ở hướng khác.
\end{block}

\renewcommand{\arraystretch}{1.12}
\resizebox{\linewidth}{!}{%
\begin{tabular}{lrrr}
\toprule
Opt & dead@final & structural & flicker \\
\midrule
sgd & 11.0 & 7.2 & 3.8 \\
\bad{adam} & 62.7 & \bad{55.1} & 7.6 \\
adahessian & 11.2 & 6.2 & 4.9 \\
\bad{sophiah} & 88.6 & \bad{78.3} & 10.3 \\
sassha & 9.4 & 5.8 & 3.7 \\
kfac & 9.8 & 5.8 & 4.0 \\
\bottomrule
\end{tabular}}
\end{columns}

\begin{alertblock}{Lõi chết sau cùng}
Flicker chỉ là lề mỏng 2--10\%; Adam và SophiaH vẫn giết một lõi capacity cấu trúc rất lớn ở mô hình cuối. Vì vậy ``Adam không chết vĩnh viễn theo thời gian'' không mâu thuẫn với structural dead 55\%.
\end{alertblock}
\end{frame}

\begin{frame}{Recovery: data-driven hay weight-driven?}
\begin{columns}[T]
\column{0.55\textwidth}
\centering
\includegraphics[width=\linewidth,height=0.68\textheight,keepaspectratio]{F18_data_vs_weight_recovery.png}
\column{0.42\textwidth}
\small
\renewcommand{\arraystretch}{1.12}
\resizebox{\linewidth}{!}{%
\begin{tabular}{lrrr}
\toprule
Opt & total rev. & data & \% data-driven \\
\midrule
kfac & 8.7 & 8.6 & \good{99} \\
sgd & 8.5 & 8.0 & 94 \\
adahessian & 8.6 & 7.7 & 90 \\
sgd-mom & 6.4 & 5.4 & 84 \\
sassha & 12.9 & 10.7 & 82 \\
sophiah & 2.1 & 1.4 & 67 \\
\bad{adam} & 13.6 & 0.8 & \bad{6} \\
\bottomrule
\end{tabular}}

\begin{block}{Diễn giải}
Đa số unit ``hồi sinh'' vì dữ liệu xoay qua phía khác của hyperplane, không phải vì trọng số dead học lại.
\end{block}
\end{columns}

\begin{alertblock}{Ngoại lệ Adam}
Adam là ngoại lệ weight-driven: churn do chính preconditioner adaptive khuấy trọng số, khác nhóm còn lại. Với đa số optimizer, recovery chủ yếu là data-driven flicker khi probe xoay.
\end{alertblock}
\end{frame}

\begin{frame}{Frozen-weight trapping: cần bỏ nhiễu weight decay}
\centering
\small
\renewcommand{\arraystretch}{1.16}
\resizebox{0.86\linewidth}{!}{%
\begin{tabular}{lccc}
\toprule
Opt & Ratio thô $\|\Delta w_{\rm dead}\|/\|\Delta w_{\rm alive}\|$ & Ratio tiếp tuyến sạch wd & \% chuyển động dead là radial decay \\
\midrule
kfac & 0.005 & 0.005 & 5 \\
sgd & 0.708 & \bad{0.027} & \bad{99} \\
sgd\_momentum & 0.831 & \bad{0.028} & \bad{100} \\
sophiah & 0.032 & 0.029 & 20 \\
adahessian & 0.467 & 0.082 & 88 \\
adam & 0.189 & 0.123 & 70 \\
sassha & 0.246 & 0.153 & 49 \\
\bottomrule
\end{tabular}}

\vspace{0.6em}
\begin{alertblock}{Gradient thật gần như không tới dead units}
Sau khi bỏ thành phần co radial do weight decay, dead unit nhận rất ít gradient ở gần như mọi optimizer. Trapping invariant-submanifold là cơ chế vững, nhưng riêng nó chưa tạo accuracy-LoP trong regime này.
\end{alertblock}
\end{frame}

\begin{frame}{Sharpness theo cả quá trình: trend quan trọng hơn mức tuyệt đối}
\begin{columns}[T]
\column{0.56\textwidth}
\centering
\includegraphics[width=\linewidth,height=0.67\textheight,keepaspectratio]{F20_sharpness_process.png}
\column{0.41\textwidth}
\small
\renewcommand{\arraystretch}{1.1}
\resizebox{\linewidth}{!}{%
\begin{tabular}{lcc}
\toprule
Opt & $\lambda_1^{\max}$ đầu $\to$ cuối & Trend \\
\midrule
sgd & 111.7 $\to$ 19.0 & phẳng dần \\
sgd-mom & 55.3 $\to$ 2.9 & phẳng dần \\
adam & 14.9 $\to$ 3.0 & phẳng dần \\
adahessian & 117.7 $\to$ 15.6 & phẳng dần \\
\bad{sophiah} & \bad{7.7 $\to$ 40.5} & \bad{sắc dần} \\
sassha & 2.2 $\to$ 0.29 & phẳng nhất \\
kfac & 191.4 $\to$ 88.6 & cao nhưng giảm \\
\bottomrule
\end{tabular}}

\begin{alertblock}{Trend mới là tín hiệu}
KFAC giữ sharpness cao nhưng không collapse; SophiaH là ca duy nhất sắc dần dài hạn và cũng collapse nặng nhất. Mức sharpness một điểm không đủ đọc cơ chế.
\end{alertblock}
\end{columns}
\end{frame}

\begin{frame}{Collapse theo độ sâu \& vai trò BatchNorm}
\begin{columns}[T]
\column{0.50\textwidth}
\centering
\includegraphics[width=\linewidth,height=0.66\textheight,keepaspectratio]{F11_depth_and_decomp.png}
\column{0.48\textwidth}
\centering
\includegraphics[width=\linewidth,height=0.66\textheight,keepaspectratio]{F11b_depth_BN1.png}
\end{columns}

\begin{alertblock}{BN cứu nhiều thứ, trừ Adam}
Death route cộng dồn ở layer sâu. BatchNorm cứu rank+dormancy cho hầu hết optimizer, nhưng \bad{Adam $\times$ BatchNorm} là ca bệnh lý riêng: dormant vẫn leo rất cao.
\end{alertblock}
\end{frame}

\section{Plasticity \& giới hạn}

\begin{frame}{Regime này chưa tạo accuracy-LoP: rank không dự báo acc}
\centering
\includegraphics[height=0.68\textheight,keepaspectratio]{F9_rank_vs_plasticity.png}

\begin{alertblock}{Scatter chính}
Các cell trải từ rank giảm mạnh đến rank tăng, nhưng $\Delta$test\_acc(late--early) vẫn $\geq 0$. Trong regime xoay mượt này, rank-collapse chưa chuyển thành loss of plasticity ở mức accuracy.
\end{alertblock}
\end{frame}

\begin{frame}{Giới hạn cần giữ khi diễn giải}
\begin{columns}[T]
\column{0.49\textwidth}
\begin{alertblock}{Thống kê}
\begin{itemize}\small
\item Chỉ \hl{n=1, seed 42}; chưa có error band.
\item Regime drift mượt không tạo accuracy-level LoP.
\item Cần thêm abrupt regime 16$\times$90$^\circ$ và benchmark cao chiều.
\end{itemize}
\end{alertblock}

\column{0.49\textwidth}
\begin{alertblock}{Đo lường}
\begin{itemize}\small
\item Tái dựng per-unit chỉ chắc cho BN=0; BN=1 thiếu running stats trong \texttt{theta\_points}.
\item Effective rank có trần $\approx$64 do probe batch, init thấp vì input 2-D.
\item KFAC \texttt{weight\_decay=0} là confound cho phân tích weight movement.
\end{itemize}
\end{alertblock}
\end{columns}

\begin{exampleblock}{Nguyên tắc báo cáo}
Chỉ nói \textbf{tín hiệu cơ chế sơ bộ}; tránh phát biểu ``optimizer X gây LoP'' từ dữ liệu two-moons này.
\end{exampleblock}
\end{frame}

\section{Tổng kết \& hướng tiếp}

\begin{frame}{H$'$: giả thuyết đề xuất để kiểm chứng}
\begin{alertblock}{H$'$ --- trạng thái: tín hiệu 1 seed, chưa xác nhận}
\textit{Optimizer adaptive theo tọa độ có clamp/saturation} như \bad{Adam, SophiaH} có thể dễ rơi vào bẫy \textbf{frozen-unit + rank-collapse} hơn các optimizer \textit{curvature-aware không-chéo} như \good{KFAC, SASSHA, AdaHessian}. Trục này có thể trực giao với phân đôi 1st/2nd.
\end{alertblock}

\centering
\small
\renewcommand{\arraystretch}{1.12}
\resizebox{0.9\linewidth}{!}{%
\begin{tabular}{lcccc}
\toprule
Opt & Route & Structural dead & Recovery mode & Sharpness trend \\
\midrule
sgd & diversifying & 7.2 & data-driven & phẳng dần \\
sgd-mom & cloning & 7.4 & data-driven & phẳng dần \\
\bad{adam} & \bad{death} & \bad{55.1} & \bad{weight-driven churn} & phẳng dần \\
adahessian & diversifying & 6.2 & data-driven & phẳng dần \\
\bad{sophiah} & \bad{death} & \bad{78.3} & mixed/weak recovery & \bad{sắc dần} \\
\good{sassha} & diversifying & 5.8 & data-driven & phẳng nhất \\
\good{kfac} & diversifying & 5.8 & data-driven & cao nhưng giảm \\
\bottomrule
\end{tabular}}

\end{frame}

\begin{frame}{Việc cần làm tiếp}
\begin{enumerate}\small
\item \textbf{Scale seed}: chạy $\geq$5 seed cho two-moons để có error band và kiểm tra độ ổn định của H$'$.
\item \textbf{Test trên nhiều setting/dataset khác}: abrupt two-moons 16$\times$90$^\circ$, permuted-MNIST, class-incremental/CIFAR để có accuracy collapse thật.
\end{enumerate}

\vspace{0.5em}
\begin{exampleblock}{Câu trả lời hiện tại}
Optimizer có dấu hiệu ảnh hưởng mạnh đến \textbf{cơ chế} LoP, nhưng trên two-moons drift mượt chưa đủ để trả lời chắc về \textbf{mức độ thuyên giảm plasticity loss}.
\end{exampleblock}
\end{frame}

\begin{frame}[standout]
Cảm ơn --- rất mong nhận được góp ý.
\end{frame}

\end{document}
